% Pollis drill-down: Compose > Compilation Lab, Julia beside compiled code: the shared example of every page, the
% pages one by one (Mix, Shared Library, Executable, WebAssembly, Translate to Rust, Julia Functions in DuckDB), where the programs
% go, what each page needs, and how the pages are made (14 November 2026). For users who want the detail, and for
% whoever maintains Pollis.
% Compile with: latexmk -pdf workflow_compilation_lab.tex
% Shared look for the one-page Pollis workflow handouts (same fonts and colours as the brochure).
\documentclass[10pt,a4paper]{article}

\usepackage[a4paper,left=17mm,right=17mm,top=14mm,bottom=16mm,footskip=8mm]{geometry}
\usepackage[T1]{fontenc}
\usepackage[default,semibold]{sourcesanspro}
\usepackage{microtype}
\usepackage{graphicx}
\usepackage[table]{xcolor}
\usepackage{tikz}
\usepackage[most]{tcolorbox}
\usepackage{fontawesome5}
\usepackage{enumitem}
\usepackage{fancyhdr}
\usepackage[hidelinks]{hyperref}

\graphicspath{{images/}}

% ---------------------------------------------------------------- colours (as in the brochure)
\definecolor{navy}{HTML}{1C3A5E}
\definecolor{pblue}{HTML}{3B78B5}
\definecolor{amber}{HTML}{F2B33D}
\definecolor{deepamber}{HTML}{D9902A}
\definecolor{pale}{HTML}{EEF3F9}
\definecolor{ink}{HTML}{2A2F36}
\definecolor{logonavy}{HTML}{173D63}
\definecolor{logodark}{HTML}{C98A1E}
\definecolor{logoamber}{HTML}{F2B446}
\definecolor{logocream}{HTML}{FEEFD2}
\definecolor{soft}{HTML}{5E6773}
\color{ink}

\setlength{\parindent}{0pt}
\setlength{\parskip}{3pt}
\setlist{nosep,leftmargin=*,itemsep=2pt}

\pagestyle{fancy}
\fancyhf{}
\renewcommand{\headrulewidth}{0pt}
\fancyfoot[L]{\footnotesize\color{soft}\textbf{Pollis}\enspace Bright Insights through Exploration, Modelling, Simulation and Optimisation}
\fancyfoot[R]{\footnotesize\color{soft}github.com/saragga/pollis/releases/latest}

% The Pollis logo, redrawn so it prints sharply at any size
\newcommand{\pollislogo}[1]{%
  \begin{tikzpicture}[x=#1/512,y=-#1/512,baseline=(current bounding box.center)]
    \fill[logonavy] (0,0) rectangle (512,512);
    \fill[logodark] (36,422) -- (104,256) -- (172,422) -- cycle;
    \fill[logoamber] (104,422) -- (203,180) -- (302,422) -- cycle;
    \fill[logocream] (203,422) -- (339,89) -- (476,422) -- cycle;
  \end{tikzpicture}}

% Top of the page: logo, name, kind of document (Workflow unless given, e.g. [Drill-Down]) and its subject;
% then the question and a one-line promise
\newcommand{\workflowhead}[4][Workflow]{%
  \thispagestyle{fancy}%
  \begin{minipage}[c]{13mm}\pollislogo{11mm}\end{minipage}%
  \begin{minipage}[c]{\dimexpr\linewidth-13mm\relax}
    {\color{navy}\Large\bfseries Pollis}\enspace{\color{soft}#1}\hfill
    {\color{deepamber}\small\bfseries\MakeUppercase{#2}}
  \end{minipage}\par\vspace{7mm}
  {\color{navy}\fontsize{22}{26}\selectfont\bfseries #3}\par\vspace{4pt}
  {\color{soft}\large #4}\par\vspace{5mm}}

% A menu path: \menu{Explore}{Access Data Libraries}... is written as \path{Explore \gt Access Data Libraries}
\newcommand{\gt}{\,{\color{deepamber}\bfseries ›}\,}
\newcommand{\menupath}[1]{{\footnotesize\bfseries\color{navy}#1}}

% A screenshot if images/<file> exists, otherwise an orange placeholder that says what to capture
\newlength{\shotheight}\setlength{\shotheight}{43mm}
\newcommand{\stepshot}[2]{%
  \IfFileExists{images/#1}{%
    \tcbox[enhanced, colback=black, colframe=black!70, boxrule=0.4pt, arc=2pt,
      left=0pt,right=0pt,top=0pt,bottom=0pt, boxsep=0pt, halign=center]{%
      \includegraphics[width=\dimexpr\linewidth-1pt\relax,height=\shotheight,keepaspectratio]{#1}}}{%
    \begin{tcolorbox}[enhanced, colback=white, colframe=orange!80, boxrule=0.6pt, arc=2pt,
      height=\shotheight, valign=center, halign=center, left=6pt, right=6pt]
      {\color{orange}\Large\faCamera}\par\smallskip
      {\color{orange}\footnotesize\bfseries Screenshot: #2}\par\smallskip
      {\color{orange!80!black}\scriptsize images/#1}
    \end{tcolorbox}}}

% One step: number, title, menu path, image file, what to capture, what it means
\newcommand{\step}[6]{%
  \begin{tcolorbox}
    {\color{deepamber}\Large\bfseries #1}\enspace{\color{navy}\large\bfseries #2}\par\vspace{1pt}
    \menupath{#3}\par\vspace{4pt}
    \stepshot{#4}{#5}\par\vspace{3pt}
    {\small #6}
  \end{tcolorbox}}

\newenvironment{steps}{%
  \begin{tcbraster}[raster columns=2, raster equal height=rows, raster column skip=4mm,
    raster row skip=4mm, enhanced, colback=pale, colframe=pale, boxrule=0pt, sharp corners,
    left=6pt,right=6pt,top=6pt,bottom=6pt, fontupper=\raggedright]}{\end{tcbraster}}

% Boxes for the bottom of the page
\newtcolorbox{pbox}[2][]{enhanced, colback=pale, colframe=amber, boxrule=0pt,
  leftrule=3pt, sharp corners, left=7pt, right=7pt, top=6pt, bottom=6pt,
  fontupper=\small\raggedright, before upper={{\color{navy}\normalsize\bfseries #2}\par\smallskip}, #1}
\newtcolorbox{nbox}[1][]{enhanced, colback=navy, colframe=navy, boxrule=0pt,
  sharp corners, left=8pt, right=8pt, top=6pt, bottom=6pt, coltext=white,
  fontupper=\small\raggedright, #1}

% A Next Steps item inside the bottom boxes
\newcommand{\nsitem}[2]{\item[\color{deepamber}\faChevronRight] \textbf{#1}: #2}

% Citations: \dcite{key}{Breiman (2001)} links to its entry \dref{key}{...} in the References, keyed as in pollis.bib
\newcommand{\dcite}[2]{\hyperlink{ref:#1}{\color{pblue}#2}}
\newcommand{\dref}[2]{\raisebox{1.1em}[0pt][0pt]{\hypertarget{ref:#1}{}}#2\par}
\newcommand{\ddoi}[1]{\href{https://doi.org/#1}{\color{pblue}doi:#1}}
\newcommand{\references}{\section*{References}\addcontentsline{toc}{section}{References}\footnotesize\raggedright\setlength{\parskip}{3pt}}

% The page of each heading, for Ask Pollis
% Pollis: the page of each heading, for Ask Pollis. A document that inputs this file, at the end of its preamble, writes
% <jobname>.labels beside its PDF, a line page|level|heading for each heading, the page counted from the first, the
% cover, as 1. compile_documents.jl copies it to ~/.pollis/documents, and ask_index.jl cuts the document into passages
% under its headings with it. \asklabel{1}{Title} marks a topic or a page, \asklabel{2}{Title} a heading inside it; the
% sections and subsections of a document that has them are marked as they go into its contents. A package vignette
% inputs it before its \documentclass (compile_vignette.jl), so the contents are taken over when the document begins,
% after the packages have taken them over. Read a second time, as by a copy of a document for Edit the LaTeX Source,
% which inputs it as resources/ask_labels when Pollis has put it in already (document_source.jl), it does nothing.
\ifdefined\asklabel\endinput\fi
\makeatletter
\newwrite\pollis@labels
\immediate\openout\pollis@labels=\jobname.labels
\newcommand{\asklabel}[2]{\write\pollis@labels{\the\ReadonlyShipoutCounter|#1|\unexpanded{#2}}}
\def\pollis@toc{toc}
\newcommand{\pollis@contentsline}[3]{%
  \def\pollis@file{#1}%
  \ifx\pollis@file\pollis@toc\asklabel{#2}{#3}\fi
  \pollis@addcontentsline{#1}{#2}{#3}}
% The headings of the New Methods style of the vignettes, \nmsection and \nmsub, which go into no contents
\def\pollis@asksection#1{\pollis@nmsection{#1}\asklabel{1}{#1}}
\def\pollis@asksub#1{\pollis@nmsub{#1}\asklabel{2}{#1}}
\AddToHook{begindocument/end}{%
  \let\pollis@addcontentsline\addcontentsline\let\addcontentsline\pollis@contentsline
  \ifdefined\nmsection\let\pollis@nmsection\nmsection\let\nmsection\pollis@asksection\fi
  \ifdefined\nmsub\let\pollis@nmsub\nmsub\let\nmsub\pollis@asksub\fi}
\makeatother



\usepackage{titlesec}
\usepackage{tabularx}
\usepackage{booktabs}
\usepackage{listings}

\fancyhf{}
\fancyfoot[L]{\footnotesize\color{soft}\textbf{Pollis}\enspace Drill-down: Compilation Lab}
\fancyfoot[R]{\footnotesize\color{soft}\thepage}

\titleformat{\section}{\color{navy}\large\bfseries}{\color{deepamber}\thesection}{0.6em}{}
\titleformat{\subsection}{\color{navy}\normalsize\bfseries}{\color{deepamber}\thesubsection}{0.5em}{}
\titlespacing*{\section}{0pt}{12pt}{4pt}
\titlespacing*{\subsection}{0pt}{8pt}{2pt}
\setlength{\parskip}{4pt}

\lstset{basicstyle=\ttfamily\footnotesize, columns=fullflexible, keepspaces=true, breaklines=true,
  backgroundcolor=\color{pale}, frame=leftline, rulecolor=\color{amber}, framerule=2pt, xleftmargin=6pt,
  framexleftmargin=6pt, aboveskip=4pt, belowskip=4pt, upquote=true, tabsize=2}

\usepackage{xurl}
\newcommand{\file}[1]{{\small\nolinkurl{#1}}}
\newcolumntype{L}{>{\raggedright\arraybackslash}X}
\newcommand{\ui}[1]{\textbf{#1}}
% Code with spaces in it, such as a command line, which \file would run together
\newcommand{\code}[1]{\texttt{\small #1}}
% A menu path at the start of a table cell
\newcommand{\cellpath}[1]{\leavevmode\menupath{#1}}

\begin{document}

\workflowhead[Drill-Down]{Compilation Lab}
  {Julia beside compiled code}
  {How Julia calls compiled code and is called by it, how a Julia function becomes a library, a program, a DuckDB
   function or WebAssembly, and how every page checks the result against Julia, pixel by pixel.}

\begin{nbox}
{\color{amber}\bfseries Drill-down.}\enspace
A look under the surface of the pages of \textbf{Compose \gt Compilation Lab}. It is for anyone who wants to know
what each page builds, where it puts it, what it needs on the computer, and why its check can demand exact
agreement; \S\ref{sec:machinery} is for whoever maintains Pollis. It describes the state of 14 November 2026.
\end{nbox}

\tableofcontents

\section{What Compilation Lab is}\label{sec:what}
Julia compiles every function before it runs it, through LLVM, the compiler framework of \dcite{LA2004a}{Lattner and
Adve (2004)}, but the result lives inside the Julia session. Compilation Lab is
about the boundary: Julia calling code that other compilers made, other programs calling Julia, and a Julia function
taken out of the session as a library, a program, a function of a database or a module a web page runs. Its pages are in two groups.

\begin{center}\small
\renewcommand{\arraystretch}{1.25}
\begin{tabularx}{\linewidth}{@{}p{0.36\linewidth}p{0.14\linewidth}L@{}}
\toprule
\textbf{Menu item} & \textbf{Style} & \textbf{What it does}\\
\midrule
\multicolumn{3}{@{}l}{\emph{Compile Julia}}\\
\cellpath{Compile Julia to a Shared Library} & Panel & One Julia function compiled into a library by three
  tools, called from a C program (\S\ref{sec:shared})\\
\cellpath{Compile Julia to an Executable} & Panel & One Julia program compiled into an executable by three tools
  (\S\ref{sec:exe})\\
\cellpath{Compile Julia to WebAssembly} & Steps & One Julia function compiled to WebAssembly in two ways, checked and
  drawn live in the page (\S\ref{sec:wasm})\\
\multicolumn{3}{@{}l}{\emph{Other Languages}}\\
\cellpath{Mix Julia with Compiled Code} & Panel & Julia calls C, Fortran and C++; C calls Julia
  (\S\ref{sec:mix})\\
\cellpath{Translate Julia to Rust} & Steps & A Julia function translated into Rust, built and compared
  (\S\ref{sec:rust})\\
\cellpath{Julia Functions in DuckDB} & Panel & A Julia function called from SQL, from the session or as an extension
  (\S\ref{sec:duckdb})\\
\bottomrule
\end{tabularx}
\end{center}
A \emph{panel} has tabs, inputs, Example Code and Next Steps, like the method panels of Pollis; a page of
\emph{steps} takes you through numbered steps, one after the other. Every page runs its code in the Julia REPL.

\section{One example on every page}\label{sec:example}
Every page computes the same thing, so that the pages can be read side by side and each result checked against the
same Julia function: the number of iterations of $z \mapsto z^2 + c$ before $|z|$ passes 2, for each pixel of a
picture, or the most iterations allowed when it never does. Two sets come from it.
\begin{itemize}
  \item \textbf{The Mandelbrot set}: $z$ starts at 0 and $c$ is the pixel, over the region from $-2.2$ to $0.8$
    across and from $-1.2$ to $1.2$ down.
  \item \textbf{The Julia set of $c$}: $z$ starts at the pixel and $c$ is fixed, $-0.8 + 0.156i$ unless you change
    it, over the region from $-1.6$ to $1.6$ across and from $-1$ to $1$ down.
\end{itemize}
The inputs are the same on every page: \ui{Example}, \ui{Pixels Across} (800), \ui{Iterations} (256), \ui{Real Part
of c} and \ui{Imaginary Part of c}. The pixels down follow from the region, with square pixels.

The sets of $z \mapsto z^2 + c$ that start at the pixel bear the name of \dcite{Jul1918a}{Julia (1918)}, who studied
the iteration of rational functions. \dcite{Man1980a}{Mandelbrot (1980)} drew the set of the values of the parameter
for the map $z \mapsto \lambda z(1 - z)$, which a change of variable turns into $z^2 + c$; the set of the $c$ whose
iteration from 0 stays bounded now bears his name.

\subsection{Why the check can be exact}
The iteration is written with real numbers, the same operations in the same order in every language:
\begin{lstlisting}
zr2 = zr * zr;  zi2 = zi * zi
zr2 + zi2 > 4.0 && return k
zi = 2.0 * zr * zi + ci;  zr = zr2 - zi2 + cr
\end{lstlisting}
Julia's \file{Float64}, C's \file{double}, Fortran's \file{real(c_double)}, Rust's \file{f64} and DuckDB's
\file{DOUBLE} are the same IEEE 754 number, and each operation rounds in the same way everywhere; \dcite{Gol1991a}{Goldberg
(1991)} explains how. Two things could
still change a result, and each page rules them out.
\begin{itemize}
  \item \textbf{Fused multiply-add.} A compiler may compute \code{a * b + c} with one rounding instead of two. C, C++
    and Fortran are built with \code{-ffp-contract=off}, which forbids it; Julia and Rust never fuse unless the code
    asks. The Next Step \ui{Fused and Not Fused} of Mix builds the code again with fusing allowed and shows the pixels
    that change.
  \item \textbf{Numbers passed as text.} A number written on a command line or in SQL is written with Julia's
    \file{repr}, the shortest decimal that reads back to the same \file{Float64}; C reads it with \file{strtod}, DuckDB
    as text cast to \file{DOUBLE}.
\end{itemize}
The count of every pixel, an \file{Int32}, must then agree exactly with Julia's, and every page prints the number of
pixels that differ, expected to be zero.

\section{Mix Julia with Compiled Code}\label{sec:mix}
Five ways across the boundary, one to a tab. In each, the compiled code fills a Julia matrix of counts through a
pointer to its first element, with nothing copied.

\begin{center}\small
\renewcommand{\arraystretch}{1.25}
\begin{tabularx}{\linewidth}{@{}p{0.17\linewidth}p{0.25\linewidth}L@{}}
\toprule
\textbf{Tab} & \textbf{Compiler} & \textbf{What crosses}\\
\midrule
Call C & \file{cc}, \file{gcc} or \file{clang} & A C function called with \file{@ccall}\\
Call Fortran & \file{gfortran} & A subroutine with \file{bind(c)}; $c$ crosses as a complex number by value\\
Call C++ & \file{c++}, \file{g++} or \file{clang++} & Code behind \file{extern "C"}; $c$ crosses as a
  \file{std::complex<double>}, by pointer\\
Julia Callback & \file{cc}, \file{gcc} or \file{clang} & The loop in C, calling a Julia function for each pixel through
  a pointer that \file{@cfunction} makes\\
Julia from C & \file{cc}, \file{gcc} or \file{clang} & A C program that starts Julia, includes a Julia file and calls
  a function of it, with the flags of \file{julia-config.jl}\\
\bottomrule
\end{tabularx}
\end{center}
Each tab builds in \file{~/.pollis/mix/<tab>}, under a name that carries a hash of the source and the flags, so a
change makes a new library rather than one the session already holds open. Next Steps: \ui{Picture of the Set},
\ui{Source and Build Command}, \ui{Julia and Compiled Times}, \ui{Fused and Not Fused}, \ui{Share of Pixels Inside}.

\section{Compile Julia to a Shared Library}\label{sec:shared}
One Julia function, \file{fill_counts}, compiled into a library by three tools: JuliaC.jl,
\dcite{JLC2026a}{JuliaLang Contributors (2026a)}, PackageCompiler.jl, \dcite{JLC2026b}{JuliaLang Contributors
(2026b)}, and StaticCompiler.jl, \dcite{Sho2026a}{Short (2026a)}. It takes only native types, returns
an \file{Int64} (the pixels inside), allocates nothing and throws no error, so that every tool can compile it.

\begin{center}\small
\renewcommand{\arraystretch}{1.25}
\begin{tabularx}{\linewidth}{@{}p{0.2\linewidth}p{0.2\linewidth}L@{}}
\toprule
\textbf{Tab} & \textbf{Runtime in the library} & \textbf{Build}\\
\midrule
JuliaC & Trimmed & \code{juliac -{}-output-lib -{}-compile-ccallable -{}-trim=safe}; Julia 1.12 or later\\
PackageCompiler & Whole, with a system image & \file{create_library} on a small package, Fractal; the calling program
  starts the runtime with \file{init_julia}\\
StaticCompiler & None & \file{compile_shlib} on the function of the session; only for code that never allocates\\
\bottomrule
\end{tabularx}
\end{center}
A small C program, \file{caller.c}, is linked against the library, with the library's folder recorded (rpath). It
fills memory of its own, writes the counts to \file{counts.bin} and prints the pixels inside; Julia reads the file back
and compares. The library is called from its own process, not from the Julia session, because a library with a Julia
runtime is made for programs that are not Julia. Next Steps: \ui{Picture of the Set}, \ui{Source and Build Command},
\ui{Library Contents}, \ui{Sizes and Build Times}, \ui{Start-Up and Call Time}, \ui{Call the Library from Python},
\ui{Call the Library from R}, \ui{Share of Pixels Inside}. Python calls the library with \file{ctypes}, of its
standard library; R, whose \code{.C} passes every argument as a pointer, through a wrapper of a few lines of C that
\code{R CMD SHLIB} builds. Both check their counts against Julia's.

\section{Compile Julia to an Executable}\label{sec:exe}
One Julia program compiled into an executable by the same three tools: juliac (\code{-{}-output-exe -{}-bundle
-{}-trim=safe}, starting at \file{@main}), PackageCompiler (\file{create_app}, starting at \file{julia_main}) and
StaticCompiler (\file{compile_executable} of \file{static_main}, written with StaticTools, \dcite{Kel...2026a}{Keller and contributors (2026)}, with no runtime).

\subsection{The program}
The program draws the picture itself, as a BMP file: black inside the set, and outside a colour from how fast each
pixel escapes. It can be started in three ways.
\begin{center}\small
\renewcommand{\arraystretch}{1.2}
\begin{tabularx}{\linewidth}{@{}p{0.44\linewidth}L@{}}
\toprule
\textbf{Command} & \textbf{What it does}\\
\midrule
\code{fractal}, or a double-click & The Mandelbrot set with the page's defaults, into \file{fractal.bmp} in the folder
  it starts in, then opened (\file{open}, \file{start} or \file{xdg-open})\\
\code{fractal picture.bmp} & The same picture into the file named, opening nothing\\
\code{fractal picture.bmp counts.bin} and nine numbers & The page's set and inputs: height, width, the corner, the
  spacing, the set, $c$ and the iterations; the counts also go to \file{counts.bin}\\
\bottomrule
\end{tabularx}
\end{center}
The tab checks the counts against Julia's and the picture byte for byte against the one Julia draws in the session.
\ui{Run the Program in Pollis} runs the program from the REPL with one argument, so its lines print in the REPL and
its picture shows in the plot pane. \ui{Start-Up and Run Time} sets each program beside \file{julia} running the same
code as a script, which has to compile it first. The other Next Steps are \ui{Picture of the Set}, \ui{Source and
Build Command}, \ui{Program Contents}, \ui{Sizes and Build Times} and \ui{Share of Pixels Inside}.

Four Next Steps build the program again with the tab's tool, each in a folder of its own, through
\code{build\_program}, which each tab defines for its own build. \ui{Julia Sets Along a Circle} draws the Julia set of
$c$ in frames as $c$ turns around a circle, multiplied at each frame by a number $w$ that Julia passes on the command
line, and times the program beside the same frames in the session. \ui{Options with Names} reads \code{--width},
\code{--iterations} and the other inputs by name, in any order: with Julia strings and \code{tryparse} for juliac and
PackageCompiler, and with \code{MallocString} and \code{argparse} of StaticTools for StaticCompiler. \ui{Use All the
Cores} shares the columns among threads with \code{Threads.@threads}, the program starting as many as
\code{JULIA\_NUM\_THREADS} asks for. \ui{Write a PNG with a Package} writes the picture with PNGFiles.jl: PackageCompiler
compiles the package whole, while juliac builds it only with \code{--trim=unsafe-warn}, as libpng calls back into Julia
through calls that \code{--trim=safe} cannot resolve. The last two are in the JuliaC and PackageCompiler tabs only, as
a program of StaticCompiler has neither threads nor the memory that PNGFiles needs.

\section{Translate Julia to Rust}\label{sec:rust}
A page of four steps that turns a Julia function into Rust, \dcite{MK2014a}{Matsakis and Klock (2014)}, called back
from Julia with \file{@ccall}. The Rust project's book, \dcite{KN2023a}{Klabnik and Nichols (2023)}, is its guide to
the language. The original
is never changed. Its example is \file{fractal_counts}, which returns the count of every pixel as a
\code{Matrix\{Int32\}}.
\begin{enumerate}
  \item \textbf{The Function.} The function to translate, by name, and its call; the example, or one of your own.
  \item \textbf{Translate.} For the example, the crate written out: \file{Cargo.toml} (edition 2024, a
    \file{cdylib}) and \file{src/lib.rs}, whose \file{fill_counts} writes into the matrix Julia allocates, and the Julia
    function \file{fractal_counts_rust} that calls it. For a function of your own, \ui{Translate with Agent} starts the
    coding agent with the skill translate-to-rust, as the task Translate Julia to Rust of Model Tasks and Routing, on
    the Strong tier by default.
  \item \textbf{Build.} \code{cargo build -{}-release} of the crate in \ui{Folder}, or the Julia file the agent wrote.
  \item \textbf{Compare.} The two versions checked to give the same result, then timed with BenchmarkTools, with the
    verdict of \file{judge}.
\end{enumerate}

\section{Julia Functions in DuckDB}\label{sec:duckdb}
The function of one pixel, \file{escape_time(zr, zi, cr, ci, maxiter)}, called from SQL. Each tab opens a DuckDB
database in memory with a table \file{pixels}, one row for each pixel and its coordinates, and one query computes
every count:
\begin{lstlisting}
SELECT i, j, escape_time(0.0::DOUBLE, 0.0::DOUBLE, x, y, 256::INTEGER) AS steps FROM pixels
\end{lstlisting}
\begin{center}\small
\renewcommand{\arraystretch}{1.25}
\begin{tabularx}{\linewidth}{@{}p{0.17\linewidth}LL@{}}
\toprule
\textbf{Tab} & \textbf{How DuckDB gets the function} & \textbf{Who can call it}\\
\midrule
Julia Function & \file{DuckDB.@create_scalar_function} and \file{DuckDB.register_scalar_function}, from the session &
  That database, while the session lasts\\
Rust Extension & A crate written with duckdb-neo,
  \dcite{DDBL2026a}{DuckDB Labs and contributors (2026)}, for DuckDB's C API version 2, built by cargo; Julia appends the 512
  bytes of metadata DuckDB checks; \code{LOAD} reads the file & Any DuckDB 2.0.0 on the same platform that allows
  unsigned extensions\\
\bottomrule
\end{tabularx}
\end{center}
DuckDB calls a scalar function once for each chunk of up to 2048 rows, not once per row, so a function from Julia or
from an extension runs at about the speed of its own language. Working on vectors of values rather than one row at a
time is the vectorised execution of MonetDB/X100, \dcite{BZN2005a}{Boncz, Żukowski and Nes (2005)}, which DuckDB
follows. An extension made this way loads in one version of
DuckDB only, the one written in its metadata, and needs a database opened with \file{allow_unsigned_extensions},
which cannot be set later. In Pollis, \code{using DuckDB} runs on the engine Pollis ships, DuckDB 2.0.0. Next Steps:
\ui{Picture of the Set}, \ui{The Function in DuckDB}, \ui{Query Plan}, \ui{Source and Build Command}, \ui{Query
Times}, \ui{Pixels Inside, Counted in SQL}.

\section{Compile Julia to WebAssembly}\label{sec:wasm}
A page of four steps that compiles \file{escape_time}, the function of one pixel, to WebAssembly, the code a web page
runs beside JavaScript, designed by engineers of the four major browser makers, \dcite{HRSTH...2017a}{Haas, Rossberg,
Schuff, Titzer, Holman, Gohman, Wagner, Zakai and Bastien (2017)}, and then runs the module in the page itself.
\begin{enumerate}
  \item \textbf{The Function.} The example in the Julia REPL, and its counts from \file{fractal_counts}.
  \item \textbf{Compile.} The module \file{fractal.wasm}, in one of the two ways below.
  \item \textbf{Check.} The page runs the module on the example and Pollis saves its counts as \file{counts.bin};
    Julia reads them back and prints the number of pixels that differ from the counts of step 1.
  \item \textbf{Live View.} The page draws the set with the module, again at each turn of the mouse wheel, which zooms
    where the pointer is, and at each drag, which moves the picture, with the time each picture took.
\end{enumerate}
\begin{center}\small
\renewcommand{\arraystretch}{1.25}
\begin{tabularx}{\linewidth}{@{}>{\raggedright\arraybackslash}p{0.2\linewidth}LL@{}}
\toprule
\textbf{Way} & \textbf{How the module is made} & \textbf{What it exports}\\
\midrule
Through Rust & \file{escape_time} translated to Rust beside \file{fill_counts}, the loop over the pixels, and built
  with \code{cargo build -{}-release -{}-target wasm32-unknown-unknown}; the code adds the target with rustup the first
  time & \file{escape_time}, \file{fill_counts}, \file{alloc_counts}, \file{free_counts} and \file{memory}: one call
  for a whole picture\\
With WebAssemblyCompiler.jl & The Julia function of step 1 itself, compiled through Binaryen by WebAssemblyCompiler.jl,
  \dcite{Sho2026b}{Short (2026b)}, in a Julia 1.10 that the code starts through juliaup, \code{julia +1.10}, since the
  package does not install on Julia 1.12 or later & \file{escape_time} alone: one call for each pixel, from a loop in
  JavaScript\\
\bottomrule
\end{tabularx}
\end{center}
WebAssembly computes with the same 64-bit numbers as Julia and has no fused multiply-add for single numbers, so both
modules give Julia's count at every pixel. The page reads no file itself: Pollis sends it the bytes of the module at
each Check and each Draw, so a new build is the one that runs. Next Step: \ui{The Same in JavaScript}, the function
and the loop over the pixels written in JavaScript, timed in the page beside each module built.

\section{Where the programs go}\label{sec:folders}
The pages that build a program, a library or an extension have a \ui{Folder} field with \ui{Choose...}. Its default is
a folder of \file{programs} in the open Pollis Workspace, else the same folder in \file{~/.pollis/programs}.

\begin{center}\small
\renewcommand{\arraystretch}{1.2}
\begin{tabularx}{\linewidth}{@{}p{0.34\linewidth}p{0.3\linewidth}L@{}}
\toprule
\textbf{Page} & \textbf{Folder} & \textbf{Inside it}\\
\midrule
Compile Julia to a Shared Library & \file{programs/shared_libraries} & \file{juliac}, \file{packagecompiler},
  \file{staticcompiler}\\
Compile Julia to an Executable & \file{programs/executables} & \file{juliac}, \file{packagecompiler},
  \file{staticcompiler}\\
Translate Julia to Rust & \file{programs/rust} & \file{fractal}, the crate\\
Julia Functions in DuckDB & \file{programs/duckdb_extensions} & \file{fractal}, the crate, and
  \file{fractal.duckdb_extension}\\
Compile Julia to WebAssembly & \file{programs/webassembly} & \file{rust}, with the crate \file{fractal} and
  \file{fractal.wasm}; \file{webassemblycompiler}, with \file{compile.jl}, its project and \file{fractal.wasm}\\
Mix Julia with Compiled Code & \file{~/.pollis/mix} & One folder for each tab; no Folder field\\
\bottomrule
\end{tabularx}
\end{center}
The build folders of Shared Library, Executable, Julia Functions in DuckDB and WebAssemblyCompiler.jl hold
\file{built_from.txt}: the Julia version, or the version and platform of DuckDB, and the sources the build was made
from. They build again only when it differs, so a second run takes seconds even where the first build took minutes;
cargo, for the other Rust crates, builds again only what changed. Each page ends by printing where its program,
library, extension or module is.

\section{What each page needs}\label{sec:needs}
\begin{center}\small
\renewcommand{\arraystretch}{1.25}
\begin{tabularx}{\linewidth}{@{}p{0.3\linewidth}L@{}}
\toprule
\textbf{Need} & \textbf{For}\\
\midrule
A C compiler: the Command Line Tools on macOS (\code{xcode-select -{}-install}), \file{gcc} on Linux, MinGW-w64 on
  Windows & Mix (C, C++, Callback, Julia from C); Shared Library (the calling program, and PackageCompiler); Executable
  (JuliaC and PackageCompiler)\\
\file{gfortran} & Mix, Call Fortran\\
Julia 1.12 or later & JuliaC, which the code installs with \file{Pkg.Apps.add("JuliaC")}\\
PackageCompiler.jl, StaticCompiler.jl, StaticTools.jl & The tabs of those tools\\
Rust 1.85 or later, with cargo, from rustup & Translate Julia to Rust; Julia Functions in DuckDB, Rust Extension,
  whose first build also downloads duckdb-neo; Compile Julia to WebAssembly, Through Rust, whose code adds the target
  \file{wasm32-unknown-unknown}\\
Julia 1.10, through juliaup (\code{juliaup add 1.10}) & Compile Julia to WebAssembly, With WebAssemblyCompiler.jl,
  which the code installs once in a project of its own\\
DuckDB.jl and Tables.jl & Julia Functions in DuckDB\\
Python 3; R, with \file{Rscript} & Shared Library, Call the Library from Python and Call the Library from R\\
A coding agent & Translate with Agent\\
\bottomrule
\end{tabularx}
\end{center}
A missing compiler stops the tab with a message that says how to install it.

\section{Under the surface}\label{sec:machinery}
\begin{itemize}
  \item \textbf{Panels from TOML.} Mix, Shared Library, Executable and Julia Functions in DuckDB are panels of the
    scaffold: each is a TOML file (\file{mixCompiled}, \file{sharedLib}, \file{compileExe}, \file{duckdbFunctions})
    that \code{npm run gulp compile-toml} turns into the data the shared panel code draws. Their wikis are in
    \file{media/wiki/<id>}.
  \item \textbf{Two pages of steps.} Translate Julia to Rust (\file{translateRustPage.ts}) and Compile Julia to
    WebAssembly (\file{webAssemblyPage.ts}) are pages of the Time Arrow style, with their own messages, and take the
    example's Julia code from \file{fractalExample.ts}. Compare of Translate Julia to Rust shares its code with Speed Up
    a Function.
  \item \textbf{WebAssembly in a page.} The page of Compile Julia to WebAssembly is built with the option
    \file{allowWebAssembly} of \file{buildWebviewHtml}, which adds \code{\textquotesingle wasm-unsafe-eval\textquotesingle} to its
    Content-Security-Policy, so that its script can compile a module with \file{WebAssembly.instantiate}. The bytes of
    the module and the counts of Check cross between the page and Pollis as typed arrays in its messages.
  \item \textbf{The Folder field.} An input of a panel's TOML with \code{folder = "programs/executables"} gets the
    \ui{Choose...} button, which opens the folder picker. The panel's handler posts the folder of the open workspace as
    the input's default (\file{setInputDefault}), from \file{pollisWorkspaceFolderPath}, which names a folder without
    making it. The code reads it as \file{expanduser(raw"...")}.
  \item \textbf{Names shared across tabs.} Each panel keeps what its tabs built in one dictionary of the session,
    \file{builds} (Shared Library), \file{programs} (Executable) or \file{ways} (DuckDB), so that a Compare action can
    set the tabs run in the session side by side.
  \item \textbf{The agent task.} Translate with Agent is the task \file{translateRust} of \file{modelRouting.ts}, with
    the skill translate-to-rust of Pollis Agents.
\end{itemize}

\section{Where things are}

\begin{center}\small
\renewcommand{\arraystretch}{1.2}
\begin{tabularx}{\linewidth}{@{}p{0.5\linewidth}L@{}}
\toprule
\textbf{File} & \textbf{Role}\\
\midrule
\file{.../model/browser/webviews/mixCompiled.toml} & Mix Julia with Compiled Code\\
\file{.../model/browser/webviews/sharedLib.toml} & Compile Julia to a Shared Library\\
\file{.../model/browser/webviews/compileExe.toml} & Compile Julia to an Executable\\
\file{.../model/browser/webviews/duckdbFunctions.toml} & Julia Functions in DuckDB\\
\file{.../compose/browser/translateRustPage.ts} & Translate Julia to Rust\\
\file{.../compose/browser/webAssemblyPage.ts} & Compile Julia to WebAssembly\\
\file{.../compose/browser/fractalExample.ts} & The example's Julia code, for both pages of steps\\
\file{.../compose/browser/compose.contribution.ts} & The Compilation Lab menu\\
\file{.../model/browser/handlers/scaffold.handler.ts} & The Folder field's picker and default\\
\file{.../pollisWorkspace/browser/pollisWorkspace.contribution.ts} & The workspace folder \file{programs}\\
\bottomrule
\end{tabularx}
\end{center}

\references
\dref{BZN2005a}{Boncz, P.\,A., Żukowski, M. and Nes, N.\,J. (2005). MonetDB/X100: Hyper-Pipelining Query Execution. In \emph{Proceedings of the Second Biennial Conference on Innovative Data Systems Research (CIDR)}, 225--237.}
\dref{DDBL2026a}{DuckDB Labs and contributors (2026). duckdb-neo: Ergonomic wrapper for DuckDB's v2 C API (beta). \url{https://github.com/duckdb/duckdb-rs}}
\dref{Gol1991a}{Goldberg, D. (1991). What every computer scientist should know about floating-point arithmetic. \emph{ACM Computing Surveys} 23(1), 5--48. \ddoi{10.1145/103162.103163}}
\dref{HRSTH...2017a}{Haas, A., Rossberg, A., Schuff, D.\,L., Titzer, B.\,L., Holman, M., Gohman, D., Wagner, L., Zakai, A. and Bastien, J.\,F. (2017). Bringing the Web up to Speed with WebAssembly. In \emph{Proceedings of the 38th ACM SIGPLAN Conference on Programming Language Design and Implementation (PLDI 2017)}, 185--200. ACM. \ddoi{10.1145/3062341.3062363}}
\dref{Jul1918a}{Julia, G. (1918). Mémoire sur l'itération des fonctions rationnelles. \emph{Journal de Mathématiques Pures et Appliquées} 1, 47--245. \url{http://www.numdam.org/item/JMPA_1918_8_1__47_0/}}
\dref{JLC2026a}{JuliaLang Contributors (2026a). JuliaC.jl: CLI app for compiling and bundling julia binaries, especially trimmed ones. \url{https://github.com/JuliaLang/JuliaC.jl}}
\dref{JLC2026b}{JuliaLang Contributors (2026b). PackageCompiler.jl: Compile your Julia Package. \url{https://github.com/JuliaLang/PackageCompiler.jl}}
\dref{Kel...2026a}{Keller, C.\,B. and contributors (2026). StaticTools.jl: Enabling StaticCompiler.jl-based compilation of (some) Julia code to standalone native binaries by avoiding GC allocations and llvmcall-ing all the things! \url{https://github.com/brenhinkeller/StaticTools.jl}}
\dref{KN2023a}{Klabnik, S. and Nichols, C. (2023). The Rust Programming Language. No Starch Press. \url{https://doc.rust-lang.org/book/}}
\dref{LA2004a}{Lattner, C. and Adve, V. (2004). LLVM: A Compilation Framework for Lifelong Program Analysis and Transformation. In \emph{International Symposium on Code Generation and Optimization (CGO 2004)}, 75--86. IEEE. \ddoi{10.1109/CGO.2004.1281665}}
\dref{Man1980a}{Mandelbrot, B.\,B. (1980). Fractal aspects of the iteration of $z \to \lambda z(1 - z)$ for complex $\lambda$ and $z$. \emph{Annals of the New York Academy of Sciences} 357(1), 249--259. \ddoi{10.1111/j.1749-6632.1980.tb29690.x}}
\dref{MK2014a}{Matsakis, N.\,D. and Klock, F.\,S. (2014). The Rust language. \emph{ACM SIGAda Ada Letters} 34(3), 103--104. \ddoi{10.1145/2692956.2663188}}
\dref{Sho2026a}{Short, T. (2026a). StaticCompiler.jl: Compiles Julia code to a standalone library (experimental). \url{https://github.com/tshort/StaticCompiler.jl}}
\dref{Sho2026b}{Short, T. (2026b). WebAssemblyCompiler.jl: Create WebAssembly with Julia. \url{https://github.com/tshort/WebAssemblyCompiler.jl}}

\end{document}
